% TOP_PROBLEM: {"id":30007579,"problem_number":"LOCAL-30007579","title":"Congestion in matching markets","category_id":21,"category":{"id":21,"name":"economics","display_name":"Economics","description":"Open economic theory statements, published research agendas, and proposed scoped specifications; question provenance and historical priority are qualified per record.","slug":"economics","order_index":21,"created_at":"2026-10-08T00:00:00Z"},"status":"research_question","external_url":null,"legacy_ids":[],"published":false,"statement":"Determine tradeoffs between signaling rounds, interview limits, and interim blocking pairs in matching markets with heterogeneous quality and preference learning.","economics":{"original_id":68,"field":"Mechanism design and market design","kind":"design","formalization_basis":"adapted_research_question","source_status":"research_agenda","priority_status":"earliest_found","priority_note":"This is the earliest source in which the relevant agenda was verified during this audit; absolute priority has not been established.","earliest_verified_reference":{"authors":["Alvin E. Roth"],"title":"Market Design: Understanding markets well enough to fix them when they’re broken","year":2010,"url":"https://web.stanford.edu/~alroth/papers/NSF%20Grand%20Challenge.SBE%202020.%20Market%20Design.pdf","doi":null,"locator":"§1 Fundamental questions, pp.3–4","evidence_summary":"This primary white paper explicitly identifies safe simple package bidding, management of matching congestion, and full transplant-center participation as open agendas."},"current_reference":{"authors":["Maxwell Allman","Itai Ashlagi","Amin Saberi","Sophie H. Yu"],"title":"From signaling to interviews in random matching markets","year":2025,"url":"https://arxiv.org/pdf/2501.14159","doi":null,"locator":"§6 Conclusion, p.19, October8,2025 revision","evidence_summary":"After proving single-round signaling results, the conclusion identifies multiple rounds, richer heterogeneity, and alternative preference-learning models as future directions."},"formalization_note":"The multi-round target is adapted from2025 §6. Roth2010 supplies an earlier explicit congestion agenda, not priority for this precise signaling model."},"statusQualification":"Published question or research agenda verified at the cited locator. The mathematical specification is adapted for this catalogue and includes additional assumptions; source publication does not establish first-ever priority or certify this exact model remains unsolved. The multi-round target is adapted from2025 §6. Roth2010 supplies an earlier explicit congestion agenda, not priority for this precise signaling model.","statusReviewedAt":"2026-10-08"}
% FIRST_POSTED: 2026-10-08
% ENTRY_KIND: statement_only
% !TeX program = lualatex
\documentclass[11pt]{article}
\usepackage{fontspec}
\usepackage[a4paper,margin=25mm]{geometry}
\usepackage{amsmath,amssymb,mathtools}
\usepackage{xurl}
\usepackage[hidelinks]{hyperref}
\newcommand{\sref}[2]{\hyperlink{source:#1:#2}{[S#2]}}
\setlength{\emergencystretch}{3em}
\begin{document}
% problemId:30007579
\section*{Congestion in matching markets}\label{30007579}
\subsection{Definitions and mathematical statement}
Let \(A\) be applicants and \(J\) firms in a finite one-to-one matching market. Before interviews, agent \(i\)'s utility for potential partner \(j\) is \(v_j+B_{ij}\); interviewing reveals an additional score \(C_{ij}\). The joint distribution of these scores and public values \(v_j\) is specified. A protocol uses \(r\) signaling rounds, at most \(d\) outgoing signals per agent per round, and at most \(k\) interviews per agent. Signals and interview invitations may depend on earlier messages and revealed scores. Final matches can use only interviewed pairs. An unmatched agent has utility zero. Define interim preferences using revealed total scores for interviewed partners and conditional expected scores for others. A pair blocks if both prefer each other to their final partners under these preferences. Determine the attainable tradeoff between total interviews, signaling rounds, and
\[
 E[\#\{i:i\text{ belongs to an interim blocking pair}\}],
\]
subject to Bayesian incentive compatibility of signaling and an explicitly specified invitation rule. Derive sufficient bounds and matching impossibility examples for richer public-quality heterogeneity and multi-round learning. State how correlations and unbalanced market sizes affect guarantees. Single-round random-market results and asymptotic stability thresholds already established are maintained. The current paper explicitly identifies multiple rounds and richer heterogeneity as unresolved directions, rather than leaving all congestion-management questions open.

\par\textbf{Formulation and scope.} The multi-round target is adapted from2025 §6. Roth2010 supplies an earlier explicit congestion agenda, not priority for this precise signaling model.

\par\textbf{Open-question provenance.} An explicit question or research agenda was verified at the source locator. This does not by itself certify that every later version remains unresolved \sref{30007579}{1}.

\par\textbf{Historical priority.} This is the earliest source in which the relevant agenda was verified during this audit; absolute priority has not been established.

\subsection{Short English statement}
Determine tradeoffs between signaling rounds, interview limits, and interim blocking pairs in matching markets with heterogeneous quality and preference learning.

\subsection{Sources}
\begin{itemize}\raggedright
\item[\textnormal{[S1]}]\hypertarget{source:30007579:1}{} Alvin E. Roth. 2010. Market Design: Understanding markets well enough to fix them when they’re broken. §1 Fundamental questions, pp.3–4.
\url{https://web.stanford.edu/~alroth/papers/NSF%20Grand%20Challenge.SBE%202020.%20Market%20Design.pdf}.
{\small\textit{Earliest verified question/agenda reference; historical priority qualified below. This primary white paper explicitly identifies safe simple package bidding, management of matching congestion, and full transplant-center participation as open agendas.}}

\item[\textnormal{[S2]}]\hypertarget{source:30007579:2}{} Maxwell Allman, Itai Ashlagi, Amin Saberi, Sophie H. Yu. 2025. From signaling to interviews in random matching markets. §6 Conclusion, p.19, October8,2025 revision.
\url{https://arxiv.org/pdf/2501.14159}.
{\small\textit{Supporting or current literature; see evidence qualification. After proving single-round signaling results, the conclusion identifies multiple rounds, richer heterogeneity, and alternative preference-learning models as future directions.}}
\end{itemize}
\end{document}
